\section{The terminal maximum and the lower bound}
\label{sec:terminal-lower}

The coarse construction produces a cell with about
$2L^2-2L\log L$ completed terminal pieces.  Its mean occupation is
still below the desired level because the construction stops at depth
$L=J-h$.  We now recover the missing contribution by maximizing over
the sites inside the terminal disk.  Conditional on the entrance and
exit data, the pieces are independent, and the covariance of their
occupation field is close to a multiple of the killed Green function.
A Gaussian approximation therefore reduces the remaining estimate to
the maximum of the discrete Gaussian free field.
We first prove this conditional terminal estimate, then combine it with the
lattice root event to obtain a disk maximum with positive probability.
Independent walk segments amplify that probability to an almost-sure bound
at deterministic times.

\subsection{Conditional occupation in one terminal disk}

Use the terminal domains $D_y$ and $V_y$, the endpoint law $\Lambda_M$,
and the completed-piece occupation vectors $X_i$ defined above.  Within
this subsection, write $G=G_{D_y}$.  For an endpoint pair $\omega=(a,e)$,
set
\[
 \mu_x(\omega)=E^\omega X(x),\qquad
 q_{x,x'}(\omega)=E^\omega[X(x)X(x')].
\]
The first task is to control these conditional moments simultaneously
over the terminal disk.  This is the reason for taking a smaller
terminal logarithmic scale here than in the upper bound.

\begin{lemma}[The terminal maximum]
\label{lem:terminal-maximum-lower}
Fix $1<p<2$ and put
\[
 h=(\log J)^p,\qquad M=e^h,\qquad L=J-h.
\]
Suppose an admissible centered block word has $G_0\le m_0$ blocks and
\begin{equation}
 b=2L^2-2L\log L+O(L)
 \tag{63}\label{eq:terminal-piece-count}
\end{equation}
completed terminal pieces, with a uniform constant in the remainder.
Put $t_J=\log J/(J\log\log J)$.  Call the endpoint realization
\emph{moment regular} if, for every $x,x'\in V_y$,
\begin{equation}
 \begin{split}
 \left|\frac1b\sum_{i\le b}\mu_x(\omega_i)-\Lambda_M\mu_x\right|
 &\le 2t_J,\\
 \left|\frac1b\sum_{i\le b}q_{x,x'}(\omega_i)
                    -\Lambda_Mq_{x,x'}\right|
 &\le 2t_J\bigl(1+G(x,x')\bigr).
 \end{split}
 \tag{L1}\label{eq:terminal-moment-regularity}
\end{equation}
Uniformly over these words,
\begin{equation}
 P(\text{moment regularity fails}\mid w)
 \le CM^4e^{-cbt_J^2}=o(J^{-A})
 \tag{L2}\label{eq:terminal-moment-failure}
\end{equation}
for every fixed $A>0$.

Let $\mathcal C_y$ be the intrinsic coarse sigma-field, which records
the word and its endpoint pairs and erases the bridge interiors.
For every moment-regular realization, under the conditional law given
$\mathcal C_y$,
\begin{equation}
 \max_{x\in V_y}\sum_{i\le b}X_i(x)
 =gb+2\sqrt2\,g\sqrt b\,h
 +O_P\left(J\sqrt h+\frac{J\log J}{\log\log J}\right).
 \tag{L3}\label{eq:terminal-maximum-expansion}
\end{equation}
The remainder is uniform over such realizations.  In particular, for
every fixed $\eta>0$, the conditional probability that this maximum
is at least $u_{J,2+\eta}$ tends to one uniformly.
\end{lemma}

\begin{proof}
We first identify the conditional means and covariances.  We then
approximate the centered maximum by a Gaussian maximum and compare
the latter with the discrete Gaussian free field.  The scale
$1<p<2$ makes both the simultaneous moment screen and the final
error estimate compatible with the $J\log J$ margin.

\smallskip\noindent
\underline{\textit{Step 1. Conditional moments and their common center.}}
Fix the entrance and directed-exit pairs $(a_i,e_i)_{i\le b}$.
The strong Markov property makes the killed bridge occupation vectors
independent under this conditioning.  Ordering the two visit times
gives the exact identities
\begin{equation}
 E^{a,e}X(x)=\frac{G(a,x)H(x,e)}{H(a,e)}
 \tag{65}\label{eq:terminal-exact-mean}
\end{equation}
and
\begin{equation}
 \begin{split}
 E_a[X(x)X(x');E=e]
 &=G(a,x)G(x,x')H(x',e)\\
 &\quad+G(a,x')G(x',x)H(x,e)
 -\mathbf1_{\{x=x'\}}G(a,x)H(x,e).
 \end{split}
 \tag{66}\label{eq:terminal-exact-second-moment}
\end{equation}
The subtraction accounts for the visit counted in both orders when
$x=x'$.  Dividing \eqref{eq:terminal-exact-second-moment} by $H(a,e)$
gives $q_{x,x'}(a,e)$.

The fixed-ratio Harnack estimates yield
\begin{equation}
 \|\mu_x/C\|_\infty\le1,\qquad
 \left\|\frac{q_{x,x'}}{C(1+G(x,x'))}\right\|_\infty\le1
 \tag{67}\label{eq:terminal-normalized-moments}
\end{equation}
uniformly in $x,x'$.  The exact exit-edge averages (43a)--(43b)
identify these functions with the finite Green class (47i), so the
word concentration estimate (49) applies.  There are at most $CM^4$
functions.  Moreover, admissibility and the present scale imply
\begin{equation}
 G_0/b+\varepsilon_M+Q/M=o(t_J),\qquad \eta_w^*=O(1).
 \tag{67a}\label{eq:terminal-screen-errors}
\end{equation}
Apply (49) with $t=t_J/C_1$, where $C_1$ is a sufficiently large fixed
multiple of the normalization constant in (47i).  After restoring
the normalizations in \eqref{eq:terminal-normalized-moments}, the
tolerances are those in \eqref{eq:terminal-moment-regularity}.  Thus
\begin{equation}
 \begin{split}
 P(\text{\eqref{eq:terminal-moment-regularity} fails}\mid w)
 &\le CM^4\exp\left\{-c\frac{\log^2J}{(\log\log J)^2}\right\}\\
 &=o(J^{-A})
 \end{split}
 \tag{68}\label{eq:terminal-screen-union-bound}
\end{equation}
for every fixed $A>0$, because
$h=o(\log^2J/(\log\log J)^2)$.  The estimate is uniform over the
normalized reset kernels; these are included in the joint law and
then integrated out.  This proves \eqref{eq:terminal-moment-failure}.
Only the lower bound uses this simultaneous $M^4$-function screen.
At the upper-bound scale $h=(C_h\log J)^4$, Lemma~\ref{lem:terminal-feynman-kac}
uses the one-point exponential estimate (50) before taking its
spatial union bound.

To compute the center, multiply
\eqref{eq:terminal-exact-mean}--\eqref{eq:terminal-exact-second-moment}
by $\Lambda_M(a,e)=\lambda_A(a)H(a,e)$ and sum first over $e$.
The Green identities (44)--(45) give
\begin{equation}
 \begin{split}
 \Lambda_M\mu_x&=d_x,\\
 \Lambda_Mq_{x,x'}
 &=d_xG(x,x')+d_{x'}G(x',x)-\mathbf1_{\{x=x'\}}d_x,\\
 \sup_{x\in V_y}|d_x-g|&\le C\varepsilon_M.
 \end{split}
 \tag{L4}\label{eq:terminal-common-moments}
\end{equation}
In particular, a moment-regular realization satisfies
\begin{equation}
 \max_{x\in V_y}
 \left|\frac1b\sum_{i\le b}E^{\omega_i}X_i(x)-g\right|
 \le Ct_J+O(\varepsilon_M).
 \tag{70}\label{eq:terminal-conditional-mean}
\end{equation}
Write
\[
 \widetilde X_i(x)=X_i(x)-E^{\omega_i}X_i(x),\qquad
 U_x=b^{-1/2}\sum_{i\le b}\widetilde X_i(x).
\]
Since $\sup_{x,x'}G(x,x')\le Ch$, the second-moment error in the
screen is $O(t_J(1+h))=o(1)$.  The error from replacing $d_x$ by $g$
in \eqref{eq:terminal-common-moments} is $O(\varepsilon_Mh)=o(1)$.
The diagonal correction is bounded, as is
\[
 \frac1b\sum_{i\le b}
 E^{\omega_i}X_i(x)E^{\omega_i}X_i(x')
\]
by \eqref{eq:terminal-normalized-moments}.  Consequently
\begin{equation}
 \max_{x,x'\in V_y}
 \left|\operatorname{Cov}(U_x,U_{x'})-2gG_{D_y}(x,x')\right|
 \le C.
 \tag{74}\label{eq:terminal-covariance-comparison}
\end{equation}
The exact exit-edge averaging is what supplies the common center;
no angular interpolation of a directed-edge observable is required.

\smallskip\noindent
\underline{\textit{Step 2. Approximation of the maximum by a Gaussian maximum.}}
The number of coordinates is $P=|V_y|=e^{2h+O(1)}$, whereas there
are only $b\asymp J^2$ independent vectors.  We therefore need an
approximation whose dimension dependence is logarithmic.

We use Chernozhukov--Chetverikov--Kato
\cite[Theorem~3.1]{CCK2016}.  In the centered form needed here, let
$Z_i\in\mathbb R^P$ be independent centered vectors with finite
absolute third moments, and let $Y_i$ be independent centered
Gaussian vectors with the same covariances.  Set
\[
 Z=\max_{j\le P}b^{-1/2}\sum_{i\le b}Z_{ij},\qquad
 \widetilde Z=\max_{j\le P}b^{-1/2}\sum_{i\le b}Y_{ij}.
\]
For every Borel set $A\subset\mathbb R$ and $\delta>0$,
\begin{equation}
 P(Z\in A)\le P(\widetilde Z\in A^{C_7\delta})
 +\frac{C_8\log^2P}{\delta^3\sqrt b}
   \{L_b+M_{b,Z}(\delta)+M_{b,Y}(\delta)\},
 \tag{69}\label{eq:terminal-cck}
\end{equation}
where $A^r=\{s:\operatorname{dist}(s,A)\le r\}$, the constants
are universal, and
\[
 \begin{split}
 L_b&=\max_{j\le P}\frac1b\sum_{i\le b}E|Z_{ij}|^3,\\
 M_{b,Z}(\delta)
 &=\frac1b\sum_{i\le b}E\left[
   \max_{j\le P}|Z_{ij}|^3;
   \max_{j\le P}|Z_{ij}|>\frac{\delta\sqrt b}{\log P}\right].
 \end{split}
\]
The definition of $M_{b,Y}$ is the same with $Y$ in place of $Z$.
We apply this statement conditionally, with $Z_i=\widetilde X_i$.

The rank-one identity (42) gives the needed tail estimate with a
factor $h^{-1}$.  Indeed, at $\theta=c_0/h$, with $c_0>0$ fixed
and sufficiently small, (39) gives $r_\theta\le C/h$, and hence
\[
 E^\omega(e^{\theta X(x)}-1)
 =r_\theta\mu_x(\omega)\le C/h.
\]
For $s\ge h$, exponential Markov's inequality gives
$P^\omega(X(x)>s)\le Ch^{-1}e^{-cs/h}$.  To see the same scale
for smaller $s$, the generating function in (42) identifies
$X(x)$ as zero with probability $1-\mu_x(\omega)/G(x,x)$ and,
otherwise, as a geometric random variable of mean $G(x,x)$.
Since $G(x,x)\asymp h$ on $V_y$, we obtain, uniformly in the
endpoint pair,
\begin{equation}
 P^\omega(X(x)>s)\le Ch^{-1}e^{-cs/h},\qquad
 E^\omega|X(x)-E^\omega X(x)|^3\le Ch^2.
 \tag{71}\label{eq:terminal-bridge-tails}
\end{equation}
Here and below $s\ge0$.  Integrating the tail proves the moment
bound, and therefore $L_b=O(h^2)$.

For $\delta=1$, the truncation threshold is
\[
 T=\frac{\sqrt b}{\log P}\asymp\frac Jh.
\]
A union bound over $P$ coordinates, followed by integration of
\eqref{eq:terminal-bridge-tails}, gives the truncated third moment. To
make the truncation explicit, put
$V_i=\max_{x\in V_y}|\widetilde X_i(x)|$. Since $P\le Ce^{2h}$ and
$T\asymp J/h\to\infty$, \eqref{eq:terminal-bridge-tails} implies
$P(V_i>t)\le CP h^{-1}e^{-ct/h}$ for $t\ge T$. Therefore
\[
\begin{aligned}
 E[V_i^3;V_i>T]
 &=T^3P(V_i>T)+3\int_T^\infty t^2P(V_i>t)\,dt\\
 &\le CP h^{-1}(T^3+hT^2+h^2T+h^3)e^{-cT/h}.
\end{aligned}
\]
Here $T/h\asymp J/h^2$ and $\log P=O(h)$, which gives
\begin{equation}
 \frac1b\sum_{i\le b}E^{\omega_i}\left[
   \max_{x\in V_y}|\widetilde X_i(x)|^3;
   \max_{x\in V_y}|\widetilde X_i(x)|>T\right]
 \le \exp\{-cJ/h^2+Ch\}.
 \tag{72}\label{eq:terminal-truncated-third-moment}
\end{equation}
The polynomial factors are absorbed by the negative exponent. The
matching Gaussian vectors have coordinate variances
at most $Ch$, by the second derivative of (42).  Their contribution
is bounded by $\exp\{-cJ^2/h^3+Ch\}$ after the same union bound.
Thus the error in \eqref{eq:terminal-cck} is at most
\begin{equation}
 O(h^4/J)+\exp\{-cJ/h^2+Ch\}=o(1).
 \tag{73}\label{eq:terminal-cck-error}
\end{equation}
All estimates depend on the endpoint data only through the uniform
bounds already established.

\smallskip\noindent
\underline{\textit{Step 3. The Gaussian maximum.}}
Let $\Gamma=(\Gamma_x:x\in V_y)$ be centered Gaussian with the
covariance of $U$, and let $\phi$ be the zero-boundary discrete
Gaussian free field in $D_y$, so that
\[
 E[\phi_x\phi_{x'}]=G_{D_y}(x,x').
\]
The extremal-process result of Biskup--Louidor
\cite[(1.2), (1.6), and Theorem~2.1]{BiskupLouidor2018}, together
with its stated consequence for maxima on open subsets following
(1.8), gives tightness of the maximum on a fixed inner subdisk after
centering by
\[
 m_M=2\sqrt g\log M-\frac34\sqrt g\log\log M.
\]
The digital disks here approximate a disk, and $V_y$ is a fixed
nonempty subdisk strictly inside it.  In particular,
\begin{equation}
 \max_{x\in V_y}\phi_x=2\sqrt g\,h+O_P(\log h).
 \tag{75}\label{eq:terminal-dgff-maximum}
\end{equation}
Only this localization, rather than the full limiting point process,
is used below.

To pass to expectations, convergence in distribution alone is
insufficient.  The tightness just stated places every median of the
maximum at $m_M+O(1)$.  Its maximal coordinate variance is $O(h)$,
so Gaussian concentration places the mean within $O(\sqrt h)$ of
a median.  Therefore
\begin{equation}
 E\max_{x\in V_y}\phi_x=2\sqrt g\,h+O(\sqrt h).
 \tag{75a}\label{eq:terminal-dgff-mean}
\end{equation}

We next check that the bounded covariance error in
\eqref{eq:terminal-covariance-comparison} changes the expected
maximum by at most $O(\sqrt h)$.  For $\lambda>0$, define
\[
 F_\lambda(z)=\lambda^{-1}\log\sum_{x\in V_y}e^{\lambda z_x}.
\]
If $p_x=e^{\lambda z_x}/\sum_{x'}e^{\lambda z_{x'}}$, then
\[
 \partial_{xx'}F_\lambda
 =\lambda(p_x\mathbf1_{\{x=x'\}}-p_xp_{x'}),\qquad
 \sum_{x,x'}|\partial_{xx'}F_\lambda|\le2\lambda.
\]
Gaussian interpolation between $\Gamma$ and $\sqrt{2g}\,\phi$
therefore bounds the difference of their expected $F_\lambda$
values by $C\lambda$, using
\eqref{eq:terminal-covariance-comparison}.  Since
\[
 0\le F_\lambda(z)-\max_xz_x\le\frac{\log P}{\lambda},
\]
choosing $\lambda\asymp\sqrt{\log P}$ yields
\begin{equation}
 \left|E\max_x\Gamma_x-E\max_x\sqrt{2g}\,\phi_x\right|
 \le C\sqrt{\log P}.
 \tag{76}\label{eq:terminal-gaussian-mean-comparison}
\end{equation}
Both fields have maximal coordinate variance $O(h)$.  Combining
Gaussian concentration with
\eqref{eq:terminal-dgff-mean}--\eqref{eq:terminal-gaussian-mean-comparison}
gives
\begin{equation}
 \max_x\Gamma_x=2\sqrt2\,g h+O_P(\sqrt h+\log h)
             =2\sqrt2\,g h+O_P(\sqrt h).
 \tag{77}\label{eq:terminal-gaussian-maximum}
\end{equation}

Apply \eqref{eq:terminal-cck} to each of the two half-lines.  For
any $K>0$, the probability that $\max_xU_x$ is below
$2\sqrt2\,gh-K\sqrt h$ is bounded by the corresponding Gaussian
lower-tail probability with its threshold increased by $C_7$, plus
the error in \eqref{eq:terminal-cck-error}.  The upper tail is bounded
by the Gaussian upper tail with its threshold decreased by $C_7$,
plus the same error.  Equation~\eqref{eq:terminal-gaussian-maximum}
thus transfers to $\max_xU_x$, uniformly in the regular endpoint
data.  This argument compares distributions; it requires no
pathwise coupling of the maxima.

\smallskip\noindent
\underline{\textit{Step 4. Restoring the means and the terminal depth.}}
Equation~\eqref{eq:terminal-conditional-mean} implies
\[
 \sum_{i\le b}E^{\omega_i}X_i(x)
 =gb+O\left(\frac{J\log J}{\log\log J}\right)
\]
uniformly in $x$.  Multiplying the centered maximum estimate by
$\sqrt b=O(J)$ proves \eqref{eq:terminal-maximum-expansion}.

Finally, \eqref{eq:terminal-piece-count} gives
\begin{equation}
 \begin{split}
 gb+2\sqrt2\,g\sqrt b\,h
 &=2gJ^2-2gJ\log J+O(h^2+h\log J+J).
 \end{split}
 \tag{73}\label{eq:terminal-restored-depth}
\end{equation}
For example, this follows by first using
$\sqrt b=\sqrt2L-(\log L)/\sqrt2+O(1+(\log L)^2/L)$ and
then substituting $L=J-h$.  The $4gLh$ term from the Gaussian
maximum restores the missing terminal depth in $2gL^2$.
Since $p<2$,
\[
 J\sqrt h+\frac{J\log J}{\log\log J}
 +h^2+h\log J+J=o(J\log J).
\]
The level in \eqref{eq:terminal-restored-depth} exceeds
$u_{J,2+\eta}$ by $\eta gJ\log J+o(J\log J)$, which proves
the final assertion.
\end{proof}

\subsection{A positive probability of a large disk maximum}
\label{sec:disk-lower}

The terminal estimate is conditional on the coarse data of one
fixed cell.  To combine it with the existence of a coarse root, we
sum the failures over all candidates.  Selecting a successful coarse
root first would introduce additional conditioning on bridge
interiors from other cells, which the conditional estimate does not
cover.

\begin{proposition}[The disk maximum]
\label{prop:disk-maximum-lower}
For every fixed $\eta>0$, there is $p_\eta>0$ such that, for all
sufficiently large $J$,
\begin{equation}
 P\left(M_{\tau_{e^J}}\ge u_{J,2+\eta}\right)\ge p_\eta.
 \tag{74}\label{eq:disk-maximum-positive-probability}
\end{equation}
\end{proposition}

\begin{proof}
Take $h=(\log J)^p$, $1<p<2$, and $L=J-h$.  Apply the lower-root
construction (32)--(38) at depth $L$, with terminal physical radius
$M=e^h$.  Restrict the centers to a fixed compact subdisk so that
their enlarged terminal domains are contained in $D(0,e^J)$.
This keeps a fixed positive fraction of the available area, and the
estimates are uniform on the restricted set.  The endpoint band in
(37) is
\[
 \sqrt{2T_L}\in[\rho_LL+z_0,\rho_LL+z_0+C_*],
\]
which gives the completed-piece count
\eqref{eq:terminal-piece-count}.

Lemma~\ref{lem:buffered-transfer} transfers the probability of
existence of a coarse root directly from Brownian motion; the pair
estimate is used on the Brownian side.  At each root, the terminal
pieces are ordered separately within the $m_0$ source intervals.
The inter-source connectors are exposed and omitted from the
occupation sums.  Their occupation is nonnegative, so omitting it
preserves a lower bound for the full-walk maximum.
For the lattice coarse-root count $Z_L^S$, the count assertion of
that lemma gives
\begin{equation}
 P(Z_L^S>0)\ge p_0,\qquad EZ_L^S\le C.
 \tag{78}\label{eq:disk-root-count}
\end{equation}

We must also show that the moment screen does not remove every
root.  Let $\widehat Z_L^S$ count the same radial barriers and
endpoint bands with all angular screens dropped.  The scalar
radial-word comparison (39d) and Rosen's Appendix Theorem~9.1(a)
give
\begin{equation}
 P(\widehat I_y^{S,\circ})\le Ce^{-2L},\qquad
 E\widehat Z_L^S\le C,
 \tag{79a}\label{eq:disk-unscreened-count}
\end{equation}
because the deterministic net has $O(e^{2L})$ points.  Conditional
on each complete radial word counted by $\widehat Z_L^S$,
\eqref{eq:terminal-moment-failure} bounds failure of moment
regularity by the same $\epsilon_J=o(J^{-A})$.  Sum first over
words and then over centers to obtain
\begin{equation}
 E\sum_y\mathbf1_{I_y^{S,\circ}}
       \mathbf1_{\{\text{the endpoints at }y\text{ are not moment regular}\}}
 \le\epsilon_J E\widehat Z_L^S=o(1).
 \tag{79}\label{eq:disk-irregular-roots}
\end{equation}
Using the unscreened count here is essential: the estimate for the
endpoints is conditional on a radial word, so its failure probability
must be integrated against the radial-word mass.  Equations
\eqref{eq:disk-root-count} and \eqref{eq:disk-irregular-roots}
show that a moment-regular lattice root exists with probability
at least $p_0-o(1)$.

For a fixed root $y$, erase its completed pieces from the path and
let $\mathcal C_y$ record the exterior path and every entrance and
directed-exit pair.  The event $G_y$ that $y$ is a moment-regular
coarse root is $\mathcal C_y$-measurable.  Given this sigma-field,
the erased pieces have the independent killed bridge laws used in
Lemma~\ref{lem:terminal-maximum-lower}.  Let
\[
 A_y=\left\{\max_{x\in V_y}\sum_{i\le b}X_i(x)
                       \ge u_{J,2+\eta}\right\}.
\]
The uniform conclusion of that lemma gives a deterministic
$\widetilde\epsilon_J\downarrow0$ such that
\begin{equation}
 P(G_y\cap A_y^c)
 =E[\mathbf1_{G_y}P(A_y^c\mid\mathcal C_y)]
 \le\widetilde\epsilon_JP(G_y).
 \tag{80}\label{eq:disk-fixed-root-failure}
\end{equation}
Consequently
\begin{equation}
 \begin{split}
 P(\exists y:G_y\cap A_y)
 &\ge P(\exists y:G_y)-\sum_yP(G_y\cap A_y^c)\\
 &\ge p_0-o(1)-\widetilde\epsilon_J EZ_L^S
 \ge p_0/2
 \end{split}
 \tag{81}\label{eq:disk-successful-root}
\end{equation}
for all sufficiently large $J$.  Each summand uses the intrinsic
bridge disintegration for its own cell; the union bound does not
require disjoint terminal disks.  On the event in
\eqref{eq:disk-successful-root}, the full local-time maximum is at
least the terminal maximum of a successful cell.  This proves
\eqref{eq:disk-maximum-positive-probability}.
\end{proof}

\subsection{An almost-sure lower bound at deterministic times}
\label{sec:deterministic-lower}

The disk estimate has a probability bounded away from zero.  We
amplify it using independent segments of increments.  The number of
segments may grow as a small power of $\log N$; shortening each
segment then costs only a controlled amount in the second-order
term of the maximum.

\begin{proposition}[The deterministic-time maximum]
\label{prop:deterministic-maximum-lower}
Let $N_j=\lfloor e^{2j}\rfloor$.  For every $\eta>0$, almost surely,
\[
 M_{N_j}\ge u_{j,2+\eta}
\]
for all sufficiently large $j$.
\end{proposition}

\begin{proof}
Fix $0<\delta<\eta$, and then choose $a>0$ so that
$\delta+2a<\eta$.  Proposition~\ref{prop:disk-maximum-lower}
gives a disk-success probability $p_\delta>0$.  By
Lemma~\ref{lem:elementary-walk-estimates}, choose a fixed $C$
large enough that
\[
 \sup_R P_0(\tau_R>CR^2)\le p_\delta/2.
\]
It follows by translation invariance that a walk segment of length
$CR^2$, from any starting site, has within-segment maximum local
time at least $u_{\log R,2+\delta}$ with probability at least
$p=p_\delta/2$, uniformly for all sufficiently large $R$.

Partition the first $N_j$ increments into
\[
 q_j=\lfloor(\log N_j)^a\rfloor
\]
disjoint blocks of length $m_j=\lfloor N_j/q_j\rfloor$, ignoring
the remainder.  Use the sites immediately after the increments in
each block to define its within-block local times.  The translated
block paths are independent, and each block maximum is bounded
above by the full local-time maximum.  Bounded endpoint changes
are absorbed by the strict margin in $\delta+2a<\eta$.
Choose $R_j$ by
\[
 \log R_j=\frac12\log(m_j/C)
         =j-\frac a2\log j+O(1).
\]
The loss in logarithmic radius is therefore $(a/2)\log j+O(1)$.
Substituting this into the quadratic leading term gives
\[
 2g(\log R_j)^2
 =2gj^2-2agj\log j+O(j+(\log j)^2),
\]
while
$g(\log R_j)\log\log R_j=gj\log j+O(j+(\log j)^2)$.
Hence
\begin{equation}
 u_{\log R_j,2+\delta}
 =u_{j,\,2+\delta+2a+o(1)}
 \ge u_{j,2+\eta}
 \tag{78}\label{eq:deterministic-block-level}
\end{equation}
for all sufficiently large $j$.  Independence of the blocks now
implies
\begin{equation}
 P(M_{N_j}<u_{j,2+\eta})
 \le(1-p)^{q_j}\le e^{-cj^a}.
 \tag{79}\label{eq:deterministic-failure-summable}
\end{equation}
The right-hand side is summable, so Borel--Cantelli proves the
assertion.
\end{proof}
